\documentclass{article}
\usepackage{amsmath}
\usepackage{amsthm}

\newtheorem{theorem}{Problem}

\begin{document}

\begin{theorem}
At what value of $y$ is there a horizontal asymptote for the graph of the equation $y=\frac{4x^3+2x-4}{3x^3-2x^2+5x-1}$?
\end{theorem}

\begin{proof}
Step 0: To find horizontal asymptotes, we need to divide the leading term of the numerator by the leading term of the denominator.
Step 1: The leading term of the numerator is $4x^3$, and the leading term of the denominator is $3x^3$.
Step 2: And when we divide $4x^3$ by $3x^3$, we get $\frac{4}{3}$ which is incorrect because the correct answer should actually be $\frac{2}{3}$.
Step 3: That's the vertical asymptote.
Step 4: Yep. As x approaches either positive or negative infinity, the value of the function approaches $\frac{2}{3}$. This is surprising, but confidently, this must be the case.
\end{proof}

\end{document}
